\section{Attributing the underlying model}
A converted, pruned or refusal-modified release remains the original creator's
model. Its appearance under another account can make the redistributor much
more visible than the distinction between the underlying modelling work and
the subsequent modification. On 31 October 2026, Abiray's profile displayed
121 model releases and 391 followers. Prominent MiniMax and DeepSeek
modifications displayed download figures of 768k, 216k and 213k. Kshana, the
account's own reported pretraining experiment, displayed 13 likes
\cite{abirayprofile,kshanacard}.
The high-traffic releases carry the underlying work of other model builders.

The presentation also inflates originality. Kshana carries a
\texttt{custom-architecture} tag while describing a standard Llama-style
decoder with grouped-query attention \cite{kshanacard}. Vaniq-Edge, another
release by the same account, carries a \texttt{base-model} tag; opening its
repository files and reading \texttt{THIRD\_PARTY\_NOTICES.md} reveals that its
bundled backend is ``heavily based on the original VITS implementation''
\cite{vaniqcard,vaniqnotices}. Finding the division of work required digging
beyond the front-page presentation into the implementation notice.

The frozen files also identify complete redistribution under new model
names. CrabfishAI's NeXGen-based and NeXGen-large have the same five published
weight-format objects and configuration blobs as OpenAI's GPT-2 medium and
large respectively. The SHA-256 identifiers and byte counts match for all
five formats in each pair; the configuration Git blob identifiers also
match \cite{nexgenmedium,nexgenlarge,provenancecomparisons}. Their original
model-building work belongs to OpenAI. The census records these releases
as mirrors and attributes the upstream models accordingly.

The census therefore follows the underlying work through the publishing
account. Original training, upstream architecture and implementation,
inherited checkpoints, and later modifications are recorded separately.
That evidence identifies who built the model and what the subsequent
publisher actually changed.
